
% =============================================================================
% Project : THE PHILOSOPHY OF VIABILITY
% Book    : Book I — Ontology of Viability
% Book ID : VS.B01
% Part    : Part I — Identity and Manifestation
% Part ID : VS.B01.P01
% Chapter : Chapter 08 — Comparative Philosophical Examination
% Chapter ID : VS.B01.P01.C08
%
% File :
%   2_MAINMATTER/PART_I_IDENTITY_AND_MANIFESTATION/
%   CHAPTER_08_COMPARATIVE_PHILOSOPHICAL_EXAMINATION/
%   sections/section_07_leibniz.tex
%
% Section :
%   Section 07 — Gottfried Wilhelm Leibniz
%
% Structural Identifier :
%   VS.B01.P01.C08.S07
%
% Purpose :
%   To reconstruct Leibniz's account of individual substance,
%   the Complete Individual Concept, and the Identity of
%   Indiscernibles; examine the relation between predicate
%   completeness and ontological identity; and critically
%   assess the explanatory force and limits of Leibniz's
%   position from the commitments and distinctions established
%   in Part I of the Ontology of Viability.
%
% Scope :
%   The section concentrates on the ontology of individual
%   substance, the containment of predicates in the subject,
%   the Complete Individual Concept, the Identity of
%   Indiscernibles, the Principle of Sufficient Reason,
%   monadic substance, unity, expression, and the problem
%   of identity.
%
% Philosophical Restriction :
%   The critique is grounded only in distinctions and
%   commitments already developed in Part I. No later
%   epistemological framework or unintroduced principle
%   of the Philosophy of Viability is adopted.
%
% Status :
%   Historical reconstruction is conditionally supported
%   by primary texts and specialist scholarship.
%   The critical assessment is argued but restricted.
%   The question whether Leibniz's Complete Individual
%   Concept entails an ontological reduction of identity
%   remains open to further interpretation and argument.
%
% =============================================================================


\section{Gottfried Wilhelm Leibniz}
\label{sec:ch08-leibniz}

\index[name]{Leibniz}
\index[name]{Max Black}
\index[name]{Adams, Robert Merrihew}
\index[concept]{Identity of Indiscernibles}
\index[concept]{Indiscernibility of Identicals}
\index[concept]{Complete Individual Concept}
\index[concept]{Predicate-in-Notion Principle}
\index[concept]{Principle of Sufficient Reason}
\index[concept]{Intrinsic Properties}
\index[concept]{Extrinsic Properties}
\index[concept]{Substance}
\index[concept]{Monad}
\index[concept]{Pre-established Harmony}
\index[concept]{Numerical Identity}
\index[concept]{Unity}
\index[concept]{Expression}
\index[concept]{Manifestation}
\index[concept]{Thisness}


% =============================================================================
% SCOPE AND COMPARATIVE ORIENTATION
% =============================================================================

\subsection*{Scope and Comparative Orientation}

% Structural ID: VS.B01.P01.C08.S07.SS01.P001
\phantomsection\label{para:VS.B01.P01.C08.S07.SS01.P001}

Leibniz occupies a particularly important position in the
comparative examination of identity because he develops
a systematic connection between individual substance,
the complete concept of an individual, the containment
of predicates in the subject, and the impossibility of
two substances being entirely alike while remaining
numerically distinct. His account is therefore not
merely a theory of how individuals may be distinguished.
It is an attempt to explain why an individual substance
is the particular substance that it is and why no other
substance can be exactly alike in every relevant respect.

% Structural ID: VS.B01.P01.C08.S07.SS01.P002
\phantomsection\label{para:VS.B01.P01.C08.S07.SS01.P002}

The central comparative question is whether the complete
predicative determination of an individual substance
can provide an exhaustive account of its identity.
The question must be formulated carefully. It is not
enough to ask whether a finite investigator can discover
all the properties of an existent. Leibniz's strongest
position does not depend upon the completeness of a
human inventory of properties. It concerns the complete
notion of the individual substance and the relation
between that notion and the truth of every predicate
attributable to it. The present examination must therefore
address the ontological force of this relation rather
than rely upon the limitations of ordinary description.


% =============================================================================
% 1. HISTORICAL AND TEXTUAL LOCATION
% =============================================================================

\subsection{Historical and Textual Location}
\label{subsec:ch08-leibniz-historical}

% Structural ID: VS.B01.P01.C08.S07.SS02.P001
\phantomsection\label{para:VS.B01.P01.C08.S07.SS02.P001}

Gottfried Wilhelm Leibniz (1646--1716) developed his
metaphysical position across several decades and through
a large body of published and unpublished writings.
His account of substance, identity, determination,
possibility, and the relation between individual beings
cannot be reconstructed responsibly from one isolated
formulation. The terminology and argumentative emphasis
of the earlier and later writings are related, but
their development must not be treated as though Leibniz
had presented one finished doctrine in an entirely
unchanging form. \cite{LookLeibnizSEP2020}

% Structural ID: VS.B01.P01.C08.S07.SS02.P002
\phantomsection\label{para:VS.B01.P01.C08.S07.SS02.P002}

The principal text for the present examination is the
\textit{Discourse on Metaphysics} (1686), especially
Sections 8, 9, and 13. Section 8 develops the account
of individual substance and the complete notion of
the subject. Section 9 presents the principle that
two substances cannot be entirely alike while differing
only numerically. Section 13 addresses the relation
between the complete notion of an individual and the
distinction between necessary and contingent truths.
These sections must be read together because the
complete concept, the identity of the individual,
and the modal status of predicates are closely
connected in Leibniz's argument. \cite[§§8--9, 13]{LeibnizDiscourse}

% Structural ID: VS.B01.P01.C08.S07.SS02.P003
\phantomsection\label{para:VS.B01.P01.C08.S07.SS02.P003}

The \textit{Monadology} (1714), especially Sections
1--9, 13--18, and 56--57, presents a more concentrated
account of simple substances, their internal
differentiation, and their relation to the universe.
It is particularly important because Leibniz does not
treat the distinction between substances as a merely
external comparison. Each monad possesses its own
determinate internal state and expresses the universe
in its own manner. The relation between the monad's
individuality and its internal character therefore
provides a further formulation of the problem examined
in the \textit{Discourse}. \cite[§§1--9, 13--18, 56--57]{LeibnizMonadology}

% Structural ID: VS.B01.P01.C08.S07.SS02.P004
\phantomsection\label{para:VS.B01.P01.C08.S07.SS02.P004}

The \textit{New Essays on Human Understanding}, particularly
Book II, Chapter 27, and the correspondence with Samuel
Clarke provide additional material concerning identity,
diversity, intrinsic difference, and the relation between
the individual and its predicates. These texts are used
here as historical sources for reconstructing Leibniz's
own position. Their discussions of knowledge, personal
identity, or other subjects are not imported as
foundational commitments into Book I. \cite[II.xxvii]{LeibnizNewEssays}
\cite{LeibnizClarkeCorrespondence}


% =============================================================================
% 2. LEIBNIZ'S METAPHYSICAL ACCOUNT OF INDIVIDUAL SUBSTANCE
% =============================================================================

\subsection{Individual Substance and the Containment of Predicates}
\label{subsec:ch08-leibniz-individual-substance}

% Structural ID: VS.B01.P01.C08.S07.SS03.P001
\phantomsection\label{para:VS.B01.P01.C08.S07.SS03.P001}

Leibniz begins his account of individual substance by
examining the relation between a subject and the
predicates attributed to it. A substance is not
adequately understood merely as a subject to which
various predicates happen to be attached. Such a
description may identify the grammatical or nominal
role of a subject, but it does not yet explain the
nature of the individual substance. Leibniz therefore
asks what makes a predicate truly attributable to
a particular subject and what relation must obtain
between the notion of the subject and the predicates
that belong to it. \cite[§8]{LeibnizDiscourse}

% Structural ID: VS.B01.P01.C08.S07.SS03.P002
\phantomsection\label{para:VS.B01.P01.C08.S07.SS03.P002}

His answer is expressed through the principle commonly
called the Predicate-in-Notion Principle: in every
true predication, the predicate is contained in the
notion of the subject. The relation is not restricted
to predicates that are explicitly present in an
ordinary verbal definition. A predicate may be
implicitly contained in the subject's notion even
when its connection with the subject is not evident
to a finite investigator. The complete understanding
of the subject would disclose the ground of every
true predicate attributable to it. \cite[§8]{LeibnizDiscourse}
\cite{LookLeibnizSEP2020}

% Structural ID: VS.B01.P01.C08.S07.SS03.P003
\phantomsection\label{para:VS.B01.P01.C08.S07.SS03.P003}

This account is connected with Leibniz's conception
of truth. A proposition is not true merely because
a predicate has been associated with a subject by
an arbitrary act of description. Its truth has a
ground in the nature of the subject. The complete
notion of an individual substance is consequently
not an incomplete description to which further
predicates may be externally added without limit.
It is a notion sufficiently complete to contain
the predicates that are truly attributable to
that individual. \cite[§8]{LeibnizDiscourse}

% Structural ID: VS.B01.P01.C08.S07.SS03.P004
\phantomsection\label{para:VS.B01.P01.C08.S07.SS03.P004}

The principle of sufficient reason provides a further
dimension of Leibniz's metaphysics. A true fact must
have a sufficient reason why it is so rather than
otherwise. Leibniz connects this requirement with
the complete notion of an individual: the predicates
true of an individual are not an arbitrary collection
of facts, but belong to an ordered account of that
individual and its place within the total order of
the world. The principle is therefore relevant to
his attempt to explain why one individual is
distinguishable from another. \cite{LookLeibnizSEP2020}
\cite[§§8--9]{LeibnizDiscourse}

% Structural ID: VS.B01.P01.C08.S07.SS03.P005
\phantomsection\label{para:VS.B01.P01.C08.S07.SS03.P005}

The complete notion also has a modal dimension.
Leibniz distinguishes necessary truths, whose
opposites involve contradiction, from contingent
truths whose certainty depends upon the order of
the world and the free decrees involved in its
actualization. The fact that a predicate is contained
in the complete notion of an individual does not
therefore mean that every predicate is necessary
in the same manner as a truth of logic or geometry.
This qualification is essential: the Complete
Individual Concept is not simply a finite definition
from which every event follows through a short
deductive chain. \cite[§13]{LeibnizDiscourse}


% =============================================================================
% 3. THE IDENTITY OF INDISCERNIBLES
% =============================================================================

\subsection{The Identity of Indiscernibles}
\label{subsec:ch08-leibniz-identity-indiscernibles}

% Structural ID: VS.B01.P01.C08.S07.SS04.P001
\phantomsection\label{para:VS.B01.P01.C08.S07.SS04.P001}

In Section 9 of the \textit{Discourse on Metaphysics},
Leibniz states that two substances cannot be entirely
alike while differing only in number. The principle
is commonly known as the Identity of Indiscernibles.
Its central claim is that numerical difference between
substances cannot obtain without a corresponding
difference in their relevant qualitative or intrinsic
determinations. Leibniz does not present the principle
merely as a practical rule for distinguishing objects.
He treats it as a consequence of the metaphysical
constitution of individual substances. \cite[§9]{LeibnizDiscourse}

% Structural ID: VS.B01.P01.C08.S07.SS04.P002
\phantomsection\label{para:VS.B01.P01.C08.S07.SS04.P002}

A familiar contemporary formulation represents the
principle in the following schematic form:

\[
\forall F\in\mathcal{P}\,
\bigl[
F(a)\leftrightarrow F(b)
\bigr]
\Rightarrow a=b.
\]

Here, \(\mathcal{P}\) designates the relevant class
of properties or predicates. The formula expresses
the idea that if two entities share every property
within the specified class, they cannot be numerically
distinct. It is a useful logical representation,
but it must not be mistaken for a complete statement
of Leibniz's historical position. The domain of
entities and the relevant class of properties require
interpretation. \cite{RodriguezPereyraIdentitySEP2025}

% Structural ID: VS.B01.P01.C08.S07.SS04.P003
\phantomsection\label{para:VS.B01.P01.C08.S07.SS04.P003}

The Identity of Indiscernibles must also be distinguished
from the Indiscernibility of Identicals. The latter
states that if two references concern one and the
same entity, the predicates truly attributable to
one are also attributable to the other. It can be
represented schematically as:

\[
a=b
\Rightarrow
\forall F\,
\bigl[
F(a)\leftrightarrow F(b)
\bigr].
\]

The two principles have opposite directions of
implication. The Indiscernibility of Identicals
does not by itself establish the Identity of
Indiscernibles. The latter is the stronger and
more controversial claim because it excludes
the possibility of two numerically distinct
entities that are alike in every relevant
non-numerical respect. \cite{RodriguezPereyraIdentitySEP2025}

% Structural ID: VS.B01.P01.C08.S07.SS04.P004
\phantomsection\label{para:VS.B01.P01.C08.S07.SS04.P004}

The scope of the principle must not be generalized
without qualification. Contemporary discussions
distinguish versions that quantify over all
properties, pure properties, or intrinsic pure
properties. These formulations are not equivalent.
A principle that includes properties referring
directly to the individual whose identity is in
question may become trivially true, while a
principle restricted to qualitative or intrinsic
properties makes a more substantive claim.
The distinction is therefore not merely technical:
it determines what kind of difference is required
for two entities to be numerically distinct.
\cite{RodriguezPereyraIdentitySEP2025}

% Structural ID: VS.B01.P01.C08.S07.SS04.P005
\phantomsection\label{para:VS.B01.P01.C08.S07.SS04.P005}

Leibniz's intended position is commonly associated
with the claim that purely extrinsic differences
cannot provide the ultimate basis for distinguishing
substances. A difference in location, for example,
does not by itself establish an intrinsic difference
between two substances. Leibniz maintains that
differences that appear relational or external
must ultimately be grounded in the determinate
nature of the substances concerned. His position
therefore cannot be reduced to the claim that
any observable difference, however accidental
or external, is sufficient to establish a
difference of substance. \cite[II.xxvii]{LeibnizNewEssays}
\cite{LookLeibnizSEP2020}


% =============================================================================
% 4. THE COMPLETE INDIVIDUAL CONCEPT
% =============================================================================

\subsection{The Complete Individual Concept}
\label{subsec:ch08-leibniz-complete-concept}

% Structural ID: VS.B01.P01.C08.S07.SS05.P001
\phantomsection\label{para:VS.B01.P01.C08.S07.SS05.P001}

The Complete Individual Concept is the central
explanatory notion through which Leibniz connects
the identity of an individual substance with
the totality of predicates true of it. The concept
is complete because it contains, or permits the
deduction of, every predicate that truly belongs
to the individual, including predicates concerning
its past, present, and future. The individual
is therefore not characterized through a
limited collection of selected properties.
Its complete notion is intended to express
the full determinate content of the individual
substance. \cite[§8]{LeibnizDiscourse}

% Structural ID: VS.B01.P01.C08.S07.SS05.P002
\phantomsection\label{para:VS.B01.P01.C08.S07.SS05.P002}

This notion must not be weakened into the idea
of an ordinary description or a list compiled
by an observer. Leibniz's claim is not that
a sufficiently detailed human description
will eventually become identical with the
individual. The Complete Individual Concept
belongs to the metaphysical account of the
substance itself and, in its fullest form,
is associated with the divine understanding.
The predicates are not merely collected because
someone happens to notice them. Their truth
is grounded in the nature of the individual
substance. \cite[§8]{LeibnizDiscourse}
\cite{LookLeibnizSEP2020}

% Structural ID: VS.B01.P01.C08.S07.SS05.P003
\phantomsection\label{para:VS.B01.P01.C08.S07.SS05.P003}

Leibniz's example of Alexander the Great illustrates
the scope of this claim. The complete notion
of Alexander contains not only predicates
that characterize him at a particular moment,
but also the predicates concerning the events
of his life and the relations in which he
stands. The notion is sufficiently determinate
to distinguish this individual from an
incomplete or general concept that might
apply to more than one possible individual.
The example is intended to show that the
individual's complete notion includes the
grounds of the predicates truly attributable
to that individual. \cite[§8]{LeibnizDiscourse}

% Structural ID: VS.B01.P01.C08.S07.SS05.P004
\phantomsection\label{para:VS.B01.P01.C08.S07.SS05.P004}

The complete concept is consequently more than
a criterion for identifying an individual
from among several candidates. It is intended
to explain why the predicates belong to that
individual and why the individual has the
determinate nature that it has. The notion
is therefore assigned a metaphysical function.
It is not simply an external representation
that happens to correspond with an independently
constituted substance. This is the strongest
form of Leibniz's position and must be preserved
in any serious critical examination. \cite{LookLeibnizSEP2020}

% Structural ID: VS.B01.P01.C08.S07.SS05.P005
\phantomsection\label{para:VS.B01.P01.C08.S07.SS05.P005}

Nevertheless, a question remains concerning
the relation between the completeness of the
individual concept and the ontological status
of the identity of the individual. Leibniz's
account establishes that all true predicates
are included in the complete notion. It does
not follow without further argument that the
identity of the individual is nothing other
than the totality of those predicates. The
difference between a complete account of
what is true of an individual and an account
of what makes that individual the individual
that it is remains the central question
for the present comparison.


% =============================================================================
% 5. MONADS, UNITY, AND INTERNAL DIFFERENCE
% =============================================================================

\subsection{Monads, Unity, and Internal Difference}
\label{subsec:ch08-leibniz-monads}

% Structural ID: VS.B01.P01.C08.S07.SS06.P001
\phantomsection\label{para:VS.B01.P01.C08.S07.SS06.P001}

In the \textit{Monadology}, Leibniz develops an
account of simple substances called monads.
A monad is simple in the sense that it has
no parts and cannot be divided into smaller
substances. Composite entities are explained
through collections or aggregates of simple
substances, while the monad itself is treated
as a genuine unity. This distinction is
central to Leibniz's account of what can
qualify as an individual substance. \cite[§§1--6]{LeibnizMonadology}

% Structural ID: VS.B01.P01.C08.S07.SS06.P002
\phantomsection\label{para:VS.B01.P01.C08.S07.SS06.P002}

Leibniz also maintains that no two monads are
exactly alike. If simple substances had no
qualitative differences, there would be no
basis for distinguishing them within the
metaphysical order. The differences between
monads are therefore not merely differences
of external position or numerical designation.
Each monad possesses a determinate internal
character through which it differs from
every other monad. \cite[§§8--9]{LeibnizMonadology}

% Structural ID: VS.B01.P01.C08.S07.SS06.P003
\phantomsection\label{para:VS.B01.P01.C08.S07.SS06.P003}

The internal character of a monad involves
a succession of perceptions and appetitions.
Leibniz describes perception as a state
that involves multiplicity within unity,
while appetition concerns the internal
principle of transition from one perception
to another. This account allows Leibniz
to explain internal differentiation without
attributing spatial parts to the simple
substance. The monad remains simple in
composition while possessing a plurality
of internal states. \cite[§§13--18]{LeibnizMonadology}

% Structural ID: VS.B01.P01.C08.S07.SS06.P004
\phantomsection\label{para:VS.B01.P01.C08.S07.SS06.P004}

Each monad also expresses the universe in
its own manner. The universe is not represented
from one identical perspective in every
substance; rather, each monad expresses the
same universal order according to its own
determinate point of view. Leibniz uses this
account to explain the correspondence among
substances without treating one created
substance as directly transferring its
internal states to another. \cite[§§56--57]{LeibnizMonadology}

% Structural ID: VS.B01.P01.C08.S07.SS06.P005
\phantomsection\label{para:VS.B01.P01.C08.S07.SS06.P005}

The distinction between simple substances
and aggregates is equally important.
Leibniz does not regard every organized
or apparently unified body as a genuine
substance in the same sense as a monad.
The unity of an aggregate and the unity
of a simple substance therefore occupy
different positions in his metaphysics.
This is a substantive commitment of his
system and cannot be transferred directly
into the Ontology of Viability. It provides
a point of comparison concerning unity,
but it does not establish that the identity
of an existent is equivalent to its unity
or to its compositional structure. \cite{LookLeibnizSEP2020}


% =============================================================================
% 6. THE STRONGEST LEIBNIZIAN DEFENCE
% =============================================================================

\subsection{The Strongest Leibnizian Defence}
\label{subsec:ch08-leibniz-strongest-defence}

% Structural ID: VS.B01.P01.C08.S07.SS07.P001
\phantomsection\label{para:VS.B01.P01.C08.S07.SS07.P001}

A serious critique must confront the strongest
form of Leibniz's position. The Complete Individual
Concept is not a finite set of predicates selected
by an observer, and the Identity of Indiscernibles
is not a rule that applies only when two entities
happen to appear alike under a limited description.
Leibniz claims that every individual substance
has a complete notion containing all predicates
true of it, and that the complete notion belongs
to the metaphysical constitution of that substance.
The fact that a human investigator cannot discover
every predicate therefore does not weaken the
principle in the manner that an objection based
solely on incomplete observation would suggest.
\cite[§§8--9]{LeibnizDiscourse}
\cite{LookLeibnizSEP2020}

% Structural ID: VS.B01.P01.C08.S07.SS07.P002
\phantomsection\label{para:VS.B01.P01.C08.S07.SS07.P002}

The strongest defence also emphasizes that the
predicates contained in the complete concept
are not all necessary in the same logical
sense. Contingent predicates are included
because they are true of the individual
within the actual order of the world, not
because their opposites are contradictory
in themselves. Leibniz therefore distinguishes
the complete inclusion of predicates from
the claim that every event is logically
necessary. The Complete Individual Concept
is intended to preserve the determinate
individuality of the substance while
maintaining the distinction between
necessary and contingent truths. \cite[§13]{LeibnizDiscourse}

% Structural ID: VS.B01.P01.C08.S07.SS07.P003
\phantomsection\label{para:VS.B01.P01.C08.S07.SS07.P003}

A further defence concerns the relation between
intrinsic and extrinsic differences. Leibniz
does not accept that two substances can be
identical in their intrinsic nature and
nevertheless differ only because they occupy
different external positions. Differences
that appear relational must, on his account,
be grounded in the intrinsic determinations
of the substances concerned. The principle
therefore has a stronger metaphysical scope
than a simple comparison of observable
features. \cite[II.xxvii]{LeibnizNewEssays}
\cite{LookLeibnizSEP2020}
\cite{RodriguezPereyraIdentitySEP2025}

% Structural ID: VS.B01.P01.C08.S07.SS07.P004
\phantomsection\label{para:VS.B01.P01.C08.S07.SS07.P004}

Leibniz's appeal to the Principle of Sufficient
Reason provides another route to the principle.
In the correspondence with Clarke, he argues
that if two substances were entirely
indiscernible, the arrangement of the world
could be altered by exchanging their
positions without producing any relevant
difference. There would then be no sufficient
reason for the actual arrangement rather
than the alternative. Since the world must
have a sufficient reason for being as it is,
Leibniz rejects the possibility of two
substances that are indiscernible in the
relevant sense. \cite[L.V.21]{LeibnizClarkeCorrespondence}
\cite{RodriguezPereyraIdentitySEP2025}

% Structural ID: VS.B01.P01.C08.S07.SS07.P005
\phantomsection\label{para:VS.B01.P01.C08.S07.SS07.P005}

These arguments show why the critique cannot
simply appeal to the fact that descriptions
may be incomplete, that investigators may
overlook differences, or that different
descriptions may be formulated for different
purposes. Such observations do not reach
the strongest Leibnizian position. The
relevant question is whether the metaphysical
completeness of the individual concept
establishes the ontological reducibility
of identity to the totality of predicates,
or whether it establishes a different relation
between individual substance and the truths
that belong to it.


% =============================================================================
% 7. CRITICAL EXAMINATION FROM THE COMMITMENTS OF PART I
% =============================================================================

\subsection{Identity and the Limits of Predicative Reduction}
\label{subsec:ch08-leibniz-critical-examination}

% Structural ID: VS.B01.P01.C08.S07.SS08.P001
\phantomsection\label{para:VS.B01.P01.C08.S07.SS08.P001}

The first critical point concerns the distinction
between the ontological identity of an existent
and the determinations through which that existent
is characterized. Part I has established that
Identity cannot be treated as an additional
property attached to an otherwise identity-free
existent, but it has also established that
Identity is not reducible to Determination.
The presence of determinations is ontologically
relevant, and some determinations may participate
in the constitution of an existent. Nevertheless,
their relevance does not establish that the
totality of determinations is identical with
the identity of that existent.


% Structural ID: VS.B01.P01.C08.S07.SS08.P002
\phantomsection\label{para:VS.B01.P01.C08.S07.SS08.P002}

This distinction places a specific burden on
the Leibnizian account. Even if every true
predicate is contained in the complete notion
of an individual substance, it remains necessary
to establish why the complete predicative
content should be identified with the
ontological identity of that substance.
The inclusion of every true predicate may
provide a complete account of what can be
truly said of the individual. It does not
automatically establish that the identity
of the individual is nothing other than
that complete predicative content.

% Structural ID: VS.B01.P01.C08.S07.SS08.P003
\phantomsection\label{para:VS.B01.P01.C08.S07.SS08.P003}

The distinction is not eliminated by the
claim that the predicates are intrinsic,
necessary to the individual, or known in
their totality by the divine understanding.
Those qualifications strengthen the
Leibnizian account of the predicates and
their relation to the substance. They do
not, by themselves, demonstrate that
Identity and the totality of determinations
are conceptually or ontologically equivalent.
The argument still requires a premise
establishing that the complete predicative
nature of the individual exhausts what
it is for that individual to be the
individual that it is.

% Structural ID: VS.B01.P01.C08.S07.SS08.P004
\phantomsection\label{para:VS.B01.P01.C08.S07.SS08.P004}

The classical counterexample proposed by
Max Black, involving two qualitatively
indiscernible spheres, illustrates a difficulty
for unrestricted versions of the principle.
If two distinct spheres could possess the
same relevant qualitative properties while
remaining numerically distinct, the claim
that no two objects can share all such
properties would be challenged. This
counterexample does not by itself refute
Leibniz's more specific metaphysical
position. Leibniz can respond that the
spheres' different positions are not
merely external differences and that
their apparent indiscernibility does
not establish the absence of intrinsic
differences. The example is therefore
a test of the scope and assumptions of
the principle, not a completed refutation
of Leibniz. \cite{BlackIdentityIndiscernibles1952}
\cite{RodriguezPereyraIdentitySEP2025}

% Structural ID: VS.B01.P01.C08.S07.SS08.P005
\phantomsection\label{para:VS.B01.P01.C08.S07.SS08.P005}

The more fundamental question concerns the
relation between the complete concept and
the individual substance. If the complete
concept is understood as the complete
ontological nature of the substance, Leibniz
must explain why the totality of predicates
is sufficient to constitute or exhaust
the identity of the substance. If it is
understood as a complete account of the
truths that belong to an individual whose
identity is already determinate, then the
concept may explain the predicates without
thereby providing the ontological ground
of identity. These are different readings,
and neither should be silently substituted
for the other.

% Structural ID: VS.B01.P01.C08.S07.SS08.P006
\phantomsection\label{para:VS.B01.P01.C08.S07.SS08.P006}

The following table summarizes the principal
distinctions required for this comparison.
It does not present the Leibnizian account
as refuted. It identifies what the account
establishes within its own framework and
what further argument is required before
its claims can be treated as an exhaustive
account of Identity in the Ontology of Viability.

% =============================================================================
% Project : THE PHILOSOPHY OF VIABILITY
% Book    : Book I — Ontology of Viability
% Book ID : VS.B01
% Asset   : TL01 — Leibnizian Predicate Completeness and Ontological Identity
% Asset ID: VS.B01.P01.C08.S07.SS08.P006.TL01
% Path    : 4_ASSETS/tables/
% =============================================================================

\begin{table}[htbp]
\centering
\small
\caption{Predicate completeness and ontological identity in the Leibnizian comparison}
\label{tab:leibniz-predicate-completeness}
\renewcommand{\arraystretch}{1.2}
\begin{tabular}{|p{0.20\textwidth}|p{0.35\textwidth}|p{0.35\textwidth}|}
\hline
\textbf{Question} & \textbf{Leibnizian commitment} & \textbf{Comparative question from Part I} \\
\hline
Individual concept &
The complete concept contains every predicate that is truly attributable to the individual. &
Does the complete determination of an individual exhaustively explain its ontological identity? \\
\hline
Identity of indiscernibles &
No two distinct individual substances differ only numerically while sharing all their relevant determinations. &
Does the absence of discernible difference establish the identity of the individuals, or only constrain their possible distinction? \\
\hline
Determination and identity &
The individual is characterized through the complete set of predicates belonging to its concept. &
Part I distinguishes Identity from Determination; the relevance of determinations does not by itself establish their identity with Identity. \\
\hline
Unity and individuality &
A monad is a simple substance and a genuine unity within Leibniz's metaphysics. &
Unity and simplicity do not, without further argument, provide an exhaustive account of Identity. \\
\hline
Expression &
Each monad expresses the universe according to its own determinate manner. &
Expression must not be identified with Identity or transferred without qualification into the distinction between Expression and Manifestation in Part I. \\
\hline
Persistence and change &
The monad's succession of internal states is explained within Leibniz's account of simple substance. &
This account does not automatically settle the broader questions of persistence, continuity, and transformation raised in Part I. \\
\hline
\end{tabular}
\end{table}

% =============================================================================
% END OF ASSET TL01
% =============================================================================



% =============================================================================
% 8. FROM COMPLETE CONCEPT TO IDENTITY: THE INFERENTIAL QUESTION
% =============================================================================

\subsection{From Complete Concept to Identity: The Inferential Question}
\label{subsec:ch08-leibniz-inferential-question}

% Structural ID: VS.B01.P01.C08.S07.SS09.P001
\phantomsection\label{para:VS.B01.P01.C08.S07.SS09.P001}

The central difficulty can now be formulated more
precisely. Leibniz establishes a relation between
the complete notion of an individual substance
and all predicates that are truly attributable
to it. He also maintains that no two substances
can be entirely alike while differing only
numerically. What remains to be examined is
the precise inferential relation between
these claims. The completeness of a concept
and the numerical uniqueness of its bearer
are related, but the relation between them
requires an argument that identifies the
premises connecting predicative completeness
with the impossibility of indiscernible
substances. \cite[§§8--9]{LeibnizDiscourse}
\cite{LookLeibnizSEP2020}

% Structural ID: VS.B01.P01.C08.S07.SS09.P002
\phantomsection\label{para:VS.B01.P01.C08.S07.SS09.P002}

One possible reconstruction is that each
individual substance has a complete concept,
and that no two distinct substances can
possess the same complete concept. If the
concepts of two substances were identical,
there would be no basis within those concepts
for distinguishing the substances. Leibniz
therefore treats the uniqueness of the
complete concept as a consequence of
the metaphysical constitution of individual
substances. This reconstruction is powerful,
but it raises the question of whether the
uniqueness of the concept has been established
independently of the Identity of Indiscernibles
or whether the principle is already involved
in the claim that every complete concept
belongs to only one substance. \cite{LookLeibnizSEP2020}

% Structural ID: VS.B01.P01.C08.S07.SS09.P003
\phantomsection\label{para:VS.B01.P01.C08.S07.SS09.P003}

The difficulty becomes sharper if the complete
concept includes the individual thisness
or the very fact that the substance is
the individual that it is. If the concept
is complete because it already contains
the identity of its bearer, then the
uniqueness of the concept may depend upon
the identity it is supposed to explain.
The argument would then establish that
each individual has a complete concept
that uniquely identifies it, but it would
not yet establish that the identity of
the individual is constituted by the
totality of qualitative predicates.
The question of whether the concept
contains a primitive individual reference
or whether the reference is reducible
to qualitative content must therefore
remain explicit. \cite{AdamsPrimitiveThisness1979}
\cite{RodriguezPereyraIdentitySEP2025}

% Structural ID: VS.B01.P01.C08.S07.SS09.P004
\phantomsection\label{para:VS.B01.P01.C08.S07.SS09.P004}

This problem must not be confused with the
claim that Leibniz simply identifies identity
with an ordinary list of properties.
His complete concept is intended to be
a metaphysical notion, not a descriptive
inventory. Nevertheless, the fact that
a notion is metaphysically complete does
not, without further argument, establish
that the identity of its bearer is reducible
to the totality of predicates contained
within it. The completeness of the notion
and the ontological status of identity
remain distinguishable questions.

% Structural ID: VS.B01.P01.C08.S07.SS09.P005
\phantomsection\label{para:VS.B01.P01.C08.S07.SS09.P005}

The distinction also allows the possibility
that the Identity of Indiscernibles could
be true without the reduction of identity
to the totality of predicates. The principle
might establish that distinct substances
cannot be qualitatively indiscernible
within the relevant domain, while leaving
open the ontological question of what
identity itself is. Conversely, even if
the principle were false in some domain,
that result would not by itself establish
that identity is independent of every
determination or difference. The truth
or falsity of the principle and the
ontological reducibility of identity
are therefore related but separate
questions.

% Structural ID: VS.B01.P01.C08.S07.SS09.P006
\phantomsection\label{para:VS.B01.P01.C08.S07.SS09.P006}

The following diagram reconstructs two principal
Leibnizian routes associated with the principle:
the route from the Complete Individual Concept
and the route from the Principle of Sufficient
Reason. The diagram marks the inferential
connection that requires further defence.
It represents the structure of arguments,
not an ontological sequence or a causal
process.

% =============================================================================
% Project : THE PHILOSOPHY OF VIABILITY
% Book    : Book I — Ontology of Viability
% Book ID : VS.B01
% Asset   : DG01 — Leibnizian Routes from Complete Concept to Indiscernibles
% Asset ID: VS.B01.P01.C08.S07.SS09.P006.DG01
% Path    : 4_ASSETS/diagrams/
% =============================================================================

\begin{figure}[htbp]
\centering
\begin{tikzpicture}[
  node distance=1.25cm and 1.4cm,
  every node/.style={font=\small, align=center},
  premise/.style={draw, rounded corners, text width=0.34\textwidth,
                  minimum height=1.0cm, inner sep=6pt},
  conclusion/.style={draw, rounded corners, text width=0.34\textwidth,
                     minimum height=1.0cm, inner sep=6pt},
  question/.style={draw, dashed, rounded corners, text width=0.78\textwidth,
                   minimum height=1.0cm, inner sep=7pt},
  arr/.style={->, thick},
  openarr/.style={->, thick, dashed}
]
\node[premise] (cic) {Complete Individual Concept\\
Every true predicate belongs to the complete notion of the individual};
\node[premise, right=of cic] (psr) {Principle of Sufficient Reason\\
There must be a sufficient reason why things are as they are};
\node[conclusion, below=of cic, xshift=2.0cm] (pii) {Identity of Indiscernibles\\
No two distinct substances are wholly indiscernible};
\draw[arr] (cic.south) -- (pii.north west);
\draw[arr] (psr.south) -- (pii.north east);
\node[question, below=of pii] (gap) {Further ontological question\\
Does the completeness of an individual's predicates establish that predicate-totality
is identical with the individual's ontological Identity?};
\draw[openarr] (pii.south) -- (gap.north);
\end{tikzpicture}
\caption{Two Leibnizian routes associated with the Identity of Indiscernibles and the further question of ontological reduction. The dashed arrow marks an inference that requires additional argument; it is not presented as a conclusion established by the principle itself.}
\label{fig:leibniz-routes-to-indiscernibles}
\end{figure}

% =============================================================================
% END OF ASSET DG01
% =============================================================================



% =============================================================================
% 9. RECIPROCAL COMPARISON: IDENTITY, UNITY, EXPRESSION, AND CHANGE
% =============================================================================

\subsection{Identity, Unity, Expression, and Change}
\label{subsec:ch08-leibniz-reciprocal-comparison}

% Structural ID: VS.B01.P01.C08.S07.SS10.P001
\phantomsection\label{para:VS.B01.P01.C08.S07.SS10.P001}

The comparison with Part I first concerns the
relation between Identity and Determination.
Leibniz's complete concept gives an extensive
account of the predicates that belong to an
individual substance. This account may explain
the determinate nature of the substance and
the grounds of its true predicates. Nevertheless,
Part I distinguishes Identity from Determination
and does not treat the totality of determinations
as an exhaustive definition of identity.
The Leibnizian account therefore raises a
genuine challenge: it asks whether the
non-reduction of identity to determination
can be maintained while acknowledging that
the complete notion of an individual
contains every true predicate attributable
to it. The challenge requires an answer
from the arguments of Part I, not the
simple repetition of its foundational
distinctions.


% Structural ID: VS.B01.P01.C08.S07.SS10.P002
\phantomsection\label{para:VS.B01.P01.C08.S07.SS10.P002}

A second point concerns Unity. Leibniz's
monad is a simple substance and therefore
a genuine unity within his metaphysical
system. Yet the distinction between unity
and identity developed in Part I prevents
the unity of a substance from being treated
as a complete account of its identity.
The fact that a monad is simple and
indivisible does not by itself establish
that identity is equivalent to simplicity,
unity, or numerical singularity. Nor
does the plurality of perceptions within
a monad establish a plurality of identities.
The relation between the unity of the
substance and its identity requires an
argument distinct from the account of
its compositional simplicity. \StructuralRef{sec:identity-and-unity}
\StructuralRef{sec:identity-and-numerical-identity}

% Structural ID: VS.B01.P01.C08.S07.SS10.P003
\phantomsection\label{para:VS.B01.P01.C08.S07.SS10.P003}

Leibniz's doctrine of expression provides
another important point of comparison.
Each monad expresses the universe according
to its own determinate perspective, and
the coordination among monads is explained
through pre-established harmony. This
account is not identical with the distinction
between Expression and Manifestation
developed in Part I. The Leibnizian concept
belongs to a specific metaphysical account
of monadic representation and the order
of substances. It cannot be transferred
directly into the vocabulary of the
Ontology of Viability. Nevertheless,
it raises a relevant question about the
relation between the identity of a substance
and the determinate manner in which it
expresses the universe. \StructuralRef{sec:identity-and-signature}

% Structural ID: VS.B01.P01.C08.S07.SS10.P004
\phantomsection\label{para:VS.B01.P01.C08.S07.SS10.P004}

The account of monadic change also requires
careful comparison. Leibniz explains the
succession of internal states through
the activity of the monad and treats the
simple substance as a unity that does
not arise or perish through the ordinary
composition or division of parts. This
is a determinate answer within his own
metaphysics. It does not establish a
universal criterion of identity preservation
for every existent, nor does it settle
the open questions concerning persistence,
continuity, and transformation developed
in Part I. \cite[§§4--7, 10--15]{LeibnizMonadology}
\StructuralRef{sec:scope-and-limits-of-identity}

% Structural ID: VS.B01.P01.C08.S07.SS10.P005
\phantomsection\label{para:VS.B01.P01.C08.S07.SS10.P005}

The comparative result is therefore not
that Leibniz's account is incompatible
with every distinction established in
Part I. Rather, it is that the relation
between complete predicative determination
and ontological identity remains insufficiently
settled by the claim that every true
predicate is contained in the complete
notion of an individual substance. The
account may provide a powerful explanation
of the determinate nature of substances
while leaving open whether identity itself
is exhausted by that explanation.


% =============================================================================
% 10. COMPARATIVE VERDICT AND UNRESOLVED QUESTIONS
% =============================================================================

\subsection{Comparative Verdict and Unresolved Questions}
\label{subsec:ch08-leibniz-verdict}

% Structural ID: VS.B01.P01.C08.S07.SS11.P001
\phantomsection\label{para:VS.B01.P01.C08.S07.SS11.P001}

Leibniz's principal contribution to the
philosophical examination of identity is
his systematic attempt to connect individual
substance, the complete notion of the subject,
and the impossibility of two substances
being entirely alike while remaining
numerically distinct. His account is
stronger than a merely descriptive criterion
of identification. It treats the complete
notion as belonging to the metaphysical
constitution of the individual and connects
the distinction among substances with
the Principle of Sufficient Reason,
intrinsic difference, and the unity
of substance.

% Structural ID: VS.B01.P01.C08.S07.SS11.P002
\phantomsection\label{para:VS.B01.P01.C08.S07.SS11.P002}

The principal difficulty is the inferential
relation between the completeness of the
individual concept and the ontological
status of identity. The fact that every
true predicate belongs to the complete
notion of an individual does not, without
an additional argument, establish that
identity is identical with the totality
of those predicates. If the complete
concept already includes the individual's
thisness, its uniqueness may presuppose
the identity that the principle is intended
to explain. If the concept is instead
understood as a complete account of the
truths belonging to an already determinate
substance, its completeness does not
by itself establish the reduction of
identity to predicate-totality.

% Structural ID: VS.B01.P01.C08.S07.SS11.P003
\phantomsection\label{para:VS.B01.P01.C08.S07.SS11.P003}

The comparison therefore supports a restricted
critical thesis: Leibniz provides a powerful
account of the complete predicative nature
of individual substances, but the relation
between that completeness and the ontological
constitution of identity requires further
argument. The distinction established in
Part I between Identity and the determinations
through which an existent is characterized
places a substantive burden on any account
that treats the totality of predicates
as an exhaustive explanation of identity.
This burden is not discharged merely by
showing that no two substances share the
same complete concept. \cite{AdamsPrimitiveThisness1979}
\cite{RodriguezPereyraIdentitySEP2025}

% Structural ID: VS.B01.P01.C08.S07.SS11.P004
\phantomsection\label{para:VS.B01.P01.C08.S07.SS11.P004}

Several questions must remain open:

\begin{enumerate}
    \item Does the Complete Individual Concept
    constitute the identity of an individual
    substance, or does it provide a complete
    account of the predicates that belong
    to an individual whose identity is
    ontologically determinate?

    \item Does the uniqueness of a complete
    individual concept follow independently
    from the Identity of Indiscernibles,
    or does the argument presuppose the
    principle it seeks to establish?

    \item Are all predicates relevant to
    the principle qualitative and intrinsic,
    or does the argument depend upon
    a broader account of the relation
    between intrinsic and extrinsic
    determinations?

    \item Can the Identity of Indiscernibles
    be true without identity being reducible
    to the totality of predicates?

    \item What is the relation between the
    unity of a monad, its internal
    differentiation, and its identity?

    \item Does Leibniz's account of expression
    provide an adequate explanation of
    the relation between individual substance
    and the universe, or does it presuppose
    further commitments concerning the
    nature of individual identity?
\end{enumerate}

These questions are not conclusions about
the falsity of Leibniz's system. They identify
the premises and relations that must be
examined before a definitive judgment
can be reached.


% =============================================================================
% 11. EPISTEMIC AND COMPARATIVE STATUS
% =============================================================================

\subsection*{Epistemic and Comparative Status of the Present Section}

% Structural ID: VS.B01.P01.C08.S07.SS12.P001
\phantomsection\label{para:VS.B01.P01.C08.S07.SS12.P001}

The historical reconstruction in this section
is supported by Leibniz's primary texts and
specialist scholarship. The reconstruction
of the Complete Individual Concept, the
Predicate-in-Notion Principle, and the
Identity of Indiscernibles is intended
to represent the strongest defensible
version of Leibniz's position. It does
not imply that every interpretive question
concerning the development or scope of
his metaphysics has been settled.

% Structural ID: VS.B01.P01.C08.S07.SS12.P002
\phantomsection\label{para:VS.B01.P01.C08.S07.SS12.P002}

The functional and justificatory status
of the principal claims must be distinguished:

\begin{itemize}
    \item \textbf{Historical reconstruction:}
    conditionally supported by the cited
    primary texts and specialist scholarship.

    \item \textbf{Leibniz's Identity of
    Indiscernibles:}
    a historical metaphysical thesis
    defended by Leibniz; it is not adopted
    as a foundational commitment of
    the Ontology of Viability.

    \item \textbf{The relation between the
    Complete Individual Concept and the
    Identity of Indiscernibles:}
    a reconstructed argument whose
    inferential structure requires
    further examination.

    \item \textbf{The non-reduction challenge
    from Part I:}
    an argued comparative assessment
    grounded in the distinction between
    Identity and Determination and in
    the established principle of
    ontological non-reduction.

    \item \textbf{The possibility that
    predicate-completeness does not
    exhaust ontological identity:}
    a critical thesis supported by the
    arguments developed in this section,
    but not a completed proof against
    every interpretation of Leibniz.

    \item \textbf{The final ontological
    relation between identity and the
    totality of predicates:}
    an unresolved problem requiring
    further argument.
\end{itemize}

% Structural ID: VS.B01.P01.C08.S07.SS12.P003
\phantomsection\label{para:VS.B01.P01.C08.S07.SS12.P003}

The principal comparative thesis is an
unproved philosophical thesis and is
therefore not assigned a theorem identifier.
The registry prefix \texttt{TH} is reserved
for demonstrated theorems, while the
present assessment remains an argued
and restricted comparative conclusion.
No new axiom, definition, or foundational
commitment of the Ontology of Viability
is introduced by this section.

% Structural ID: VS.B01.P01.C08.S07.SS12.P004
\phantomsection\label{para:VS.B01.P01.C08.S07.SS12.P004}

The section does not establish that
Leibniz has definitively reduced identity
to a mere collection of descriptive
properties. Nor does it establish that
the Complete Individual Concept is
only an external representation. It
identifies the unresolved philosophical
question concerning the relation between
the complete predicative nature of an
individual and the ontological identity
of that individual. The result must
remain proportionate to the argument
actually established.


% =============================================================================
% 12. QUESTION FOR REFLECTION
% =============================================================================

\subsection*{Question for Reflection}

% Structural ID: VS.B01.P01.C08.S07.SS13.P001
\phantomsection\label{para:VS.B01.P01.C08.S07.SS13.P001}

If the complete notion of an individual
substance contains every predicate that
is truly attributable to it, does this
establish that the totality of those
predicates is the ontological identity
of the individual, or does it establish
only that the complete notion provides
an exhaustive account of what is true
of an individual whose identity still
requires an independent ontological
explanation? In answering, distinguish
the truth of the Identity of Indiscernibles
from the claim that identity is reducible
to predicate-totality.

% =============================================================================
% END OF SECTION 07 OF CHAPTER 08
% =============================================================================